% Generated draft: requires mathematical and source review.
\documentclass[11pt]{article}
\usepackage[margin=25mm]{geometry}
\usepackage{fontspec}
\setmainfont{Latin Modern Roman}
\usepackage{amsmath,amssymb,mathtools,unicode-math}
\setmathfont{Latin Modern Math}
\usepackage{enumitem,needspace,xurl,hyperref}
\setlength{\parindent}{0pt}
\setlength{\parskip}{5pt}
\setlength{\emergencystretch}{3em}
\newcommand{\N}{\mathbb{N}}
\newcommand{\Z}{\mathbb{Z}}
\newcommand{\Q}{\mathbb{Q}}
\newcommand{\R}{\mathbb{R}}
\newcommand{\C}{\mathbb{C}}
\newcommand{\F}{\mathbb{F}}
\newcommand{\A}{\mathbb{A}}
\newcommand{\PP}{\mathbb{P}}
\newcommand{\E}{\mathbb{E}}
\newcommand{\eps}{\varepsilon}
\newcommand{\dd}{\,\mathrm{d}}
\begin{document}
\section*{UnsolvedMath ID: 30004117}
\textbf{Source problem code:} catalogue entry
\subsection{Definitions and mathematical statement}
A number field is a finite extension of \(\Q\). Under the usual convention that a link complement means a complement in the three-sphere, let \(L\subset S^3\) be a finite disjoint union of smoothly embedded circles, and suppose that \(M=S^3\setminus L\) admits a complete hyperbolic metric of sectional curvature \(-1\) and finite volume. Choose an orientation and a holonomy representation \(\rho:\pi_1(M)\to\operatorname{PSL}_2(\C)\). Define the invariant trace field by
\[
k(M)=\Q\bigl(\operatorname{tr}(\widetilde{\rho(\gamma)})^2:\gamma\in\pi_1(M)\bigr)\subset\C,
\]
where \(\widetilde{\rho(\gamma)}\in\operatorname{SL}_2(\C)\) is either lift of \(\rho(\gamma)\). The squared trace is independent of the sign of the chosen lift; no compatible choice of lifts is required in this definition. We regard number fields up to \(\Q\)-isomorphism.

The classification problem, with these conventions made explicit, is to characterize the collection
\[
\left\{[K]_{\Q}: [K:\Q]<\infty,\quad \exists L\subset S^3\ \text{such that }S^3\setminus L\text{ is complete hyperbolic of finite volume and }K\cong_{\Q}k(S^3\setminus L)\right\}.
\]
Equivalently, determine necessary and sufficient conditions on a number field \(K\) for such a link \(L\) to exist.

The supplied statement does not explicitly specify the ambient three-manifold, the completeness and finite-volume conventions, or whether fields are classified as embedded subfields of \(\C\) or up to isomorphism. The formulation above uses the standard sphere-complement and finite-volume interpretation, and classification up to \(\Q\)-isomorphism; these conventions require confirmation against the original source.
\subsection{Short English statement}
Determine exactly which number fields occur as invariant trace fields of hyperbolic link complements. The displayed formulation uses complete finite-volume complements of links in \(S^3\) and identifies fields up to \(\Q\)-isomorphism; the supplied wording leaves these conventions implicit.
\subsection{Sources}
\begin{itemize}
\item Local registry: \texttt{data/unsolvedmath/problems.json}, permanent UnsolvedMath ID \(30004117\), source problem code \texttt{catalogue entry}, title \emph{Trace Fields of Hyperbolic Link Complements}. The supplied record labels the statement exact and the problem open. These are supplied registry labels, not independent verification.
\item Supplied report citation: \emph{Mini-Workshop: Reflection Groups in Negative Curvature (2020)}, Oberwolfach Reports, pp.\ 1043--1070, DOI \texttt{10.4171/owr/2019/17}. \url{https://doi.org/10.4171/owr/2019/17}. The background attributes the collection to workshop open-problem discussions. The record gives proposed year \(2019\) and names no individual proposer; the year \(2020\) in the report title is retained as supplied.
\item Supplied literature source: \emph{Invariant trace fields of link complements}. \url{https://arxiv.org/abs/1811.00679}.
\item Supplied literature source: \emph{Oberwolfach reflection-group report}. \url{https://doi.org/10.4171/owr/2019/17}.
\item Literature assessment dated exactly \texttt{2026-08-22}: the supplied curation states that infinitely many CM fields are realized by explicit families of link complements, while the general classification remains open. Here a CM field means a totally imaginary quadratic extension of a totally real number field. This assessment is supplied dated curation derived from the OWR collection, not independently verified; the record expressly requests independent verification of time-sensitive status and cited claims.
\end{itemize}
\end{document}
